\chapter{Contexto}
\label{sec:contexto}

\section{La política de fiscalización vial en la Ciudad de México}
\label{sec:contexto_politica}

En abril de 2019 la Ciudad de México sustituyó el esquema de fotomultas —una
sanción puramente económica por exceso de velocidad captada por cámara— por el
sistema Fotocívicas, en el que la infracción se traduce en la pérdida de
puntos sobre la licencia de conducir y solo de forma secundaria en una multa
monetaria. El cambio es relevante para esta tesis porque altera el incentivo
que enfrenta el conductor: \citet{quintero2023}, en un análisis de series de
tiempo interrumpido para la Ciudad de México, documentan que las políticas de
tránsito con consecuencia económica directa (2015) se asociaron a una
reducción de la mortalidad vial, mientras que el sistema de puntos sin costo
económico inmediato (2019) se asoció más bien a un incremento. Este
antecedente anticipa que el vínculo entre fiscalización y víctimas no es
mecánico, y es consistente con el resultado nulo que se documenta en el
capítulo~\ref{chap:resultados}.

La infraestructura física de fiscalización —radares fijos y móviles que
alimentan Fotocívicas— se mantuvo relativamente estable en el periodo de
estudio: pasó de 116 dispositivos en 2019 (89 fijos y 27 móviles) a 109 en
2020 y 2021, 105 en 2022 y 121 en 2023, con una composición dominada en todos
los años por los radares fijos (entre 79 y 103 unidades). El diseño empírico
de esta tesis usa como definición del \emph{tratamiento} el inventario de
2021 —91 radares fijos y 18 móviles, geolocalizados— por ser el más cercano y
predeterminado respecto al evento principal (sección~\ref{sec:eventos}).

Sobre esta infraestructura relativamente estable se superponen cambios en
\emph{quién} puede infraccionar y \emph{cómo}. En febrero de 2021 se
modificaron las facultades de infracción de los agentes de tránsito. El
cambio de mayor relevancia para esta tesis ocurrió en abril de 2022, cuando
la Secretaría de Seguridad Ciudadana puso en marcha el Programa de Renovación
de la Policía de Tránsito \citep{elfinanciero2022}. El programa sustituyó a
la generalidad de policías preventivos —que hasta entonces podían
infraccionar— por un cuerpo reducido y especializado (100 elementos en su
arranque), identificados con un brazalete y obligados a levantar cada
infracción con un dispositivo electrónico y aplicación móvil, en vez del
talonario físico. En el mismo periodo entró en vigor un nuevo reglamento de
tránsito que, en paralelo, redujo el número de oficiales autorizados para
infraccionar y acotó los motivos por los que un vehículo puede ser remitido
al corralón \citep{galvan2022}. Ambos cambios describen la misma reforma
desde ángulos distintos: menos agentes, pero con mayor capacidad de
fiscalización efectiva por agente al estar digitalizados y ser trazables.
Es esta reforma —no el reglamento de 2022 como choque homogéneo citywide— la
que se usa como evento en el diseño empírico, precisamente porque su
implementación fue gradual y territorialmente desigual, lo que genera
variación explotable entre colonias.

Un tercer momento relevante es septiembre de 2023, cuando entró en vigor una
reforma dirigida específicamente a motociclistas (uso obligatorio de
aditamentos luminosos, visor y chamarra o peto de protección), tras un
periodo de intensificación de operativos que comenzó en enero de ese año.
Esta reforma se evalúa como evento candidato en la
sección~\ref{sec:eventos} y se usa como replicación de robustez en el
capítulo~\ref{chap:resultados}, aunque finalmente no es el evento principal
por las razones que se detallan ahí.

\section{Evolución de la siniestralidad}
\label{sec:contexto_siniestralidad}

La figura~\ref{fig:serie_victimas} muestra la serie mensual de víctimas de
hechos de tránsito registradas por la Fiscalía General de Justicia de la
Ciudad de México (FGJ CDMX) entre enero de 2019 —primer mes con cobertura
sistemática en los registros— y julio de 2024. Se observan tres rasgos
relevantes para el diseño empírico. Primero, la caída de movilidad asociada a
la pandemia de COVID-19 se refleja en un mínimo de víctimas hacia mediados de
2020, seguida de una recuperación gradual conforme se normalizó la
movilidad. Segundo, 2022 fue un año particularmente severo: la ciudad
registró 719 víctimas mortales, con los motociclistas como el grupo más
afectado (39\,\% de los fallecidos), seguidos de peatones (29\,\%), pasajeros
(17\,\%), conductores (9\,\%) y ciclistas (5\,\%) \citep{navarrete2023}. Este
repunte es contemporáneo a la renovación de la Policía de Tránsito descrita
en la sección anterior, lo que motiva usar ese evento como referencia
temporal del diseño. Tercero, ni las víctimas totales ni los fallecidos
muestran una caída sostenida después de abril de 2022; si acaso, la
tendencia de mediano plazo es ligeramente ascendente hacia 2023-2024. Esta
observación anticipa, de forma puramente descriptiva, el resultado nulo que
se documenta formalmente en el capítulo~\ref{chap:resultados}.

\begin{figure}[H]
  \centering
  \includegraphics[width=\textwidth]{401_serie_victimas_2019_2024.png}
  \caption{Víctimas de hechos de tránsito en la Ciudad de México, 2019-2024.
  Las líneas verticales marcan los tres eventos candidatos discutidos en la
  sección~\ref{sec:eventos}.}
  \imagesource{elaboración propia con datos de la FGJ CDMX.}
  \label{fig:serie_victimas}
\end{figure}

\section{Distribución espacial de la fiscalización}
\label{sec:contexto_espacial}

El diseño de exposición espacial que se define en el
capítulo~\ref{chap:estrategia} descansa en una asimetría empírica entre
dónde ocurre la fiscalización y dónde ocurren las víctimas. La
figura~\ref{fig:mapa_infracciones} muestra que las infracciones de tránsito
2020-2023 se concentran de forma extrema en un número reducido de colonias
—principalmente corredores viales primarios de Miguel Hidalgo, Cuauhtémoc,
Álvaro Obregón y Venustiano Carranza—: el 10\,\% de las colonias con más
infracciones concentra el 98.8\,\% del total registrado. En contraste, la
figura~\ref{fig:mapa_victimas} muestra que las víctimas de tránsito
2019-2024 se distribuyen de forma considerablemente más homogénea en el
territorio: el mismo umbral del 10\,\% de colonias más afectadas concentra
apenas el 44.3\,\% de las víctimas.

Esta asimetría espacial es, en buena medida, lo que justifica la elección
del diseño empírico: si la fiscalización estuviera distribuida de forma tan
concentrada como las infracciones, casi toda la variación identificante
provendría de un puñado de colonias atípicas. El uso de efectos fijos de
colonia junto con una definición de tratamiento basada en buffers de
distancia a la infraestructura de radares (en vez de, por ejemplo,
infracciones per cápita) busca precisamente evitar que la estimación quede
dominada por esas colonias de muy alta concentración de multas.

\begin{figure}[H]
  \centering
  \includegraphics[width=0.85\textwidth]{402_mapa_infracciones_colonia.png}
  \caption{Concentración espacial de las infracciones de tránsito por
  colonia, 2020-2023.}
  \imagesource{elaboración propia con datos de infracciones de tránsito y
  el catálogo de colonias del IECM.}
  \label{fig:mapa_infracciones}
\end{figure}

\begin{figure}[H]
  \centering
  \includegraphics[width=0.85\textwidth]{403_mapa_victimas_colonia.png}
  \caption{Distribución espacial de las víctimas de tránsito por colonia,
  2019-jul 2024.}
  \imagesource{elaboración propia con datos de la FGJ CDMX y el catálogo de
  colonias del IECM.}
  \label{fig:mapa_victimas}
\end{figure}

\section{Eventos candidatos considerados}
\label{sec:eventos}

La línea del tiempo de política vial 2019-2024 ofrece más de un candidato
plausible para un diseño de evento. Se evaluaron tres:

\begin{itemize}
  \item \textbf{Febrero de 2021} — modificación de las facultades de
  infracción de los agentes de tránsito. Se descarta como evento principal
  porque coincide con la fase de recuperación post-COVID de la movilidad, lo
  que hace difícil separar el efecto de la reforma del efecto mecánico de
  más autos en circulación; además, la fiscalización está deprimida en ese
  periodo y el diseño no supera la prueba conjunta de tendencias previas
  (sección~\ref{sec:pretendencias}).
  \item \textbf{Abril de 2022} — renovación de la Policía de Tránsito. Es el
  único de los tres eventos con una primera etapa plausible (más
  fiscalización, aunque débil y transitoria) y con tendencias previas
  válidas para el outcome de víctimas. Se adopta como \textbf{evento
  principal} del diseño.
  \item \textbf{Septiembre de 2023} — reforma dirigida a motociclistas. Se
  descarta como evento principal por dos razones: la cobertura de
  infracciones en los datos administrativos termina en abril de 2023, por lo
  que no hay outcome de mecanismo posterior al evento; y la variación
  espacial de la exposición es más débil que en abril de 2022. Se usa, en
  cambio, como \textbf{replicación de robustez} en el
  capítulo~\ref{chap:resultados}.
\end{itemize}

El capítulo~\ref{chap:estrategia} retoma el evento de abril de 2022 y define
formalmente el tratamiento, la especificación de estudio de evento y los
supuestos de identificación del diseño.
